% !PRINTED 26
\weberpage{26}

Es möge nun $p$ aufgehen in:
\[
 a_0,a_1\ldots a_{r-1},\quad \hbox{aber nicht in }a_r,
\]
in
\[
 b_0,b_1\ldots b_{s-1},\quad \hbox{aber nicht in }b_s,
\]
dann ist nach Voraussetzung $r\leqq m$ und kann auch gleich Null sein,
wenn $p$ schon in $a_0$ nicht aufgeht. Ebenso ist $s\leqq n$.

Bilden wir nun nach (4) $c_{r+s}$ und ordnen es in folgender Weise:
\begin{equation}
\begin{aligned}
 c_{r+s}={}&a_rb_s+a_{r-1}b_{s+1}+a_{r-2}b_{s+2}+\cdots\\
            &+a_{r+1}b_{s-1}+a_{r+2}b_{s-2}+\cdots,
\end{aligned}
\tag{1}
\end{equation}
so sieht man unmittelbar, dass $c_{r+s}$ nicht durch $p$ theilbar sein
kann, wie doch angenommen war; denn das erste Glied $a_rb_s$ ist durch
$p$ nicht theilbar, während alle anderen Glieder, da sie mit einem der
Coëfficienten $a_0,a_1\ldots a_{r-1},b_0,b_1\ldots b_{s-1}$
multiplicirt sind, durch $p$ theilbar sind. Die Annahme also, dass
$C(x)$ nicht ursprünglich sei, während es $A(x)$ und $B(x)$ sind,
führt zu einem Widerspruch.

Dieser Satz lässt sich übertragen auf Functionen von mehreren
Veränderlichen. Wir nennen eine ganze rationale Function von $m$
Veränderlichen mit ganzzahligen Coëfficienten \gesperrt{ursprünglich
oder primitiv}, wenn die Coëfficienten keinen gemeinsamen Theiler
haben, und wir sprechen den Satz aus, dass \gesperrt{das Product von
zwei ursprünglichen Functionen wieder eine ursprüngliche Function
ist}.

Um seine Wahrheit einzusehen, brauchen wir nur in der oben
durchgeführten Betrachtung die Coëfficienten $a_0,a_1,a_2\ldots$,
$b_0,b_1,b_2\ldots$ nicht als ganze Zahlen, sondern als ganze rationale
Functionen von $m-1$ Veränderlichen $y,z\ldots$ anzunehmen und eine
solche Function durch eine Primzahl $p$ theilbar zu nennen, wenn alle
ihre Coëfficienten durch $p$ theilbar sind. Setzen wir dann voraus, der
zu beweisende Satz sei bereits für Functionen von $m-1$ Variablen
bewiesen, dann ergiebt die Formel (1) seine Richtigkeit für Functionen
von $m$ Variablen, also seine allgemeine Gültigkeit durch den Schluss
der vollständigen Induction oder den Schluss von $m-1$ auf $m$.

Eine imprimitive Function ist eine solche, deren ganzzahlige
Coëfficienten alle einen gemeinsamen Theiler haben. Diesen Theiler
nennen wir den \gesperrt{Theiler der Function}. Dann können wir den
bewiesenen Satz auch so aussprechen:
%CONT>>
